\section{Problem Setup and Guarantee Boundary}
\label{sec:problem}

\subsection{Data universe}

The input is a finite snapshot of OCDS releases $D$.  A deterministic data pipeline transforms
$D$ into a knowledge graph $G=(V,E)$ and a query-record projection $U_G$.  Each row
$r\in U_G$ is keyed by a contract-node identifier and exposes a controlled set of fields,
including stored buyer and supplier names, release year, procurement category and CPV code,
selected dates, a monetary value and currency when available, an additivity flag, and provenance
links.  The word \emph{record} in the remainder of the paper refers to one row of $U_G$, not to
an arbitrary text passage or to every versioned OCDS release.

The projection deliberately preserves two different notions of identity.  Graph organisation
nodes use canonical identifiers produced by entity resolution.  Query fields such as
\texttt{buyer\_name} and \texttt{supplier\_name}, however, are stored strings derived from the
first-party data.  A set of distinct stored names must therefore not be interpreted as a set of
perfectly resolved legal entities.  We retain this distinction in result descriptions and
qualitative examples.

\subsection{Questions, programs, and denotations}

Let $q$ be a natural-language question, $\mathcal{L}$ the closed typed algebra in Section
\ref{sec:algebra}, and $p\in\mathcal{L}$ a program.  Static validation assigns a root type
$\tau$ to a program,
\begin{equation}
  \Gamma \vdash p : \tau,
  \quad
  \tau\in\{\records,\values,\groups,\numbertype,\valuetype,\booltype,\ranking\},
\end{equation}
or returns a structured type error.  Deterministic evaluation gives an answer envelope
\begin{equation}
  E_G(p) = (s,\tau,y),
\end{equation}
where $s$ is an execution status and $y$ is the denotation when $s=\textsc{ok}$.  A surface
renderer may verbalise $y$, but it may not introduce a factual value not contained in $y$.

The full runtime uses a richer intermediate representation.  A structured intent $i$ describes
the requested answer signature, source constraints, derived entity sets, target constraints, and
return operation.  A compiler or conditional planner maps $i$ to a graph plan
$H=(X,A)$, where each variable $x\in X$ is a typed record or entity set and each directed edge
in $A$ denotes a dependency whose output is bound into a downstream constraint.  Compilation
produces executable low-level specifications, which are grounded and evaluated over $U_G$.

\subsection{Answerability}

Each evaluation item has one of three operational statuses:
\begin{description}
  \item[Answerable.] A program in $\mathcal{L}$ can express the requested computation and its
  implemented data preconditions are present in the frozen snapshot.
  \item[Unsupported.] The question requests an operation outside $\mathcal{L}$, such as a median
  or mean in the current benchmark grammar.  Correct behaviour is an explicit abstention.
  \item[No result.] The requested, well-formed scope has an empty denotation in $U_G$.  Depending
  on the answer signature, correct behaviour is a faithful empty/zero result or an explicit
  no-result response, not a fabricated alternative.
\end{description}
Ambiguity and insufficient evidence are additionally represented in the full runtime.  They are
not collapsed into the same state: an unsupported operator, a valid empty query, an unresolved
literal, and a missing evidence precondition require different user-facing explanations.

\subsection{Five non-equivalent notions of success}

Table \ref{tab:guarantee-boundary} defines the terminology used throughout the paper.  The
distinctions prevent a common but consequential inference: a program can be type valid and
execute successfully while still expressing the wrong question.

\begin{table}[t]
  \caption{Guarantee boundary.  Only oracle match is an offline answer-correctness measurement;
  none of the first four rows alone establishes real-world truth.}
  \label{tab:guarantee-boundary}
  \centering
  \small
  \begin{tabular}{p{0.18\linewidth}p{0.72\linewidth}}
    \toprule
    Term & Meaning in this paper \\
    \midrule
    Static validity & The object belongs to the closed grammar and all operator input/output
    types compose; depth and node limits are respected. \\
    Executability & The grounded object can be evaluated by the implemented backend without a
    structural or runtime failure. \\
    Runtime verified & The implemented compile, preflight, execution, postflight, sanity, and
    evidence conditions required for release have passed. \\
    Evidence supported & The answer is associated with the record/evidence identifiers emitted
    by execution under the implemented provenance policy. \\
    Oracle match & The typed denotation equals the frozen benchmark oracle under the stated
    exact, set, Boolean, or numeric comparison rule. \\
    Real-world truth & A claim about procurement activity outside the finite snapshot; this is
    not guaranteed by the system or benchmark. \\
    \bottomrule
  \end{tabular}
\end{table}

For an answerable compositional item, strict correctness is
\begin{equation}
  \mathbb{1}_{\mathrm{strict}} =
  \mathbb{1}[\Gamma\vdash p:\tau]
  \mathbb{1}[E_G(p).s=\textsc{ok}]
  \mathbb{1}[\operatorname{match}_{\tau}(E_G(p).y,y^*)],
  \label{eq:strict-correctness}
\end{equation}
where $y^*$ is the frozen oracle.  Set-valued answers are compared as sets of strings, rankings
preserve their defined order and tie semantics, Booleans require Boolean values on both sides,
and numeric values use the evaluator's documented tolerance.  In particular, Python's
subtyping of \texttt{bool} under \texttt{int} must not make \texttt{True} equivalent to numeric
one.  Unsupported cases are correct only when the model explicitly declines the operation;
no-result cases use their declared empty-result convention.

Runtime verification is measured independently from Equation \ref{eq:strict-correctness}.  When
both labels are available, the informative diagnostic is the $2\times2$ table of verifier
acceptance versus oracle correctness: it reveals both accepted errors and rejected correct
plans.  The present artifacts do not support every cell of this analysis for every experiment,
so we report it only where trace provenance is complete rather than inferring it from final
accuracy.
